% Results section — converted from acm_submission/paper_draft.md
% Sections 5.1-5.11 + Caveats

\section{Experimental Setup}\label{sec:setup}

Five scenarios (A--E) differ in device mix, tariff exposure and stressors:
A is the coordination baseline, B adds battery/PV assets, C adds deferrable
flexible loads, D stresses comfort bounds, E injects component failures during
the run. Every scenario is run with \texttt{rule\_based}, \texttt{multi\_agent},
\texttt{centralized}, \texttt{mpc} and \texttt{rl} controllers on seeds 42, 123,
456, 789 and 999 (125 runs per pricing regime), plus a 50-run ablation of
\texttt{llm\_agent} and \texttt{llm\_agent\_verified} against
\texttt{multi\_agent} on the same five seeds.

The controller list is the paper's comparability spine: \texttt{rule\_based} is
the non-learning reference, \texttt{multi\_agent} is the contribution,
\texttt{centralized} is a single-optimizer upper-bound-by-construction
baseline, \texttt{mpc} is a receding-horizon LP with perfect forecasts, and
\texttt{rl} is tabular Q-learning trained per scenario (60 episodes, seed 0;
policies committed under \texttt{experiments/reference/rl\_policy/}).

\section{Cost, Peak and Service}\label{sec:costpeak}

% T2 reports mean ± 95% CI per (scenario, controller) for the five core KPIs;
% T3 reports paired per-seed deltas against rule_based with Holm-adjusted
% p-values over the family of controller comparisons.

\begin{table}[t]
\centering
\caption{Paired per-seed deltas against \texttt{rule\_based}
(mean $\pm$ 95\% CI half-width; Holm within each scenario $\times$ metric
family).}
\label{tab:deltas}
\small
\begin{tabular}{llrrr}
\toprule
Scenario & Controller & $\Delta$ cost (USD) & $\Delta$ peak (kW) & $p_{\mathrm{Holm}}$ \\
\midrule
A & multi\_agent & $-15.30 \pm 0.03$ & $-20.46 \pm 2.37$ & $< 0.0001$ \\
A & mpc & $+110.98 \pm 30.56$ & $+29.21 \pm 5.39$ & $0.0011$ \\
A & centralized & $-0.07 \pm 0.47$ & $+0.46 \pm 1.68$ & $0.72$ \\
B & multi\_agent & $-43.95 \pm 3.65$ & $-34.30 \pm 4.54$ & $< 0.0001$ \\
C & multi\_agent & $-15.30 \pm 0.03$ & $-20.49 \pm 2.41$ & $< 0.0001$ \\
E & multi\_agent & $-13.26 \pm 0.03$ & $-16.74 \pm 2.46$ & $< 0.0001$ \\
\bottomrule
\end{tabular}
\end{table}

The section's spine is the \textbf{cost-versus-service Pareto front} (F3), not
a single KPI. \texttt{multi\_agent} is the cost- and peak-minimal coordinated
controller in \textbf{every} scenario. \texttt{mpc}, with perfect forecasts,
buys near-complete EV service for $+104$ to $+126$~USD and up to $+29$~kW peak.
\texttt{centralized} does not dominate: it shaves 3--4~USD off
\texttt{rule\_based} cost but its EV service is statistically indistinguishable
from \texttt{rule\_based} ($p_{\mathrm{Holm}} = 0.72$ in A/C/D).

\subsection{LMP-Regime Benchmark}\label{sec:lmp}

The M5.2 benchmark ran the full matrix under the \texttt{lmp} pricing regime
(100 runs: 4 controllers $\times$ 5 scenarios $\times$ 5 seeds) on the
aggressive-config tree. The ordering matches the pool-TOU results:
\texttt{multi\_agent} is cost-minimal in every scenario under LMP pricing
(all $p_{\mathrm{Holm}} < 0.0001$ on cost), \texttt{mpc} is the
service-maximal end of the front (EV completion 0.93--0.99), and
\texttt{centralized} is statistically indistinguishable from
\texttt{rule\_based} on cost. The LMP-regime magnitudes are slightly smaller
than the pool-TOU magnitudes because LMP prices are lower on average, but the
\textbf{ordering and significance are identical}: locational pricing does not
change the controller ranking, it changes the magnitudes.

Scenario~E under LMP reports \texttt{member\_cost\_coverage} $= 0.80$
(building\_3 interrupted) in every run, consistent with
Section~\ref{sec:scenarioe}.

\section{Scenario E: Failure Injection}\label{sec:scenarioe}

Scenario~E injects component failures mid-run. The interruption is a property
of the scenario, not of the strategy: in \textbf{all 25} Scenario~E runs, for
\textbf{every} controller, \texttt{members\_interrupted} =
\texttt{building\_3}, giving \texttt{member\_cost\_coverage} $= 0.80 \pm 0.00$.
Community cost is reported with a second coverage covariate,
\texttt{energy.cost\_coverage} $= 0.9726$ in every~E run --- identical across
controllers, so it does not confound the E comparison.

\section{Ablation: Verified vs.\ Unverified Agents}\label{sec:ablation}

All ablation runs were produced \textbf{without an LLM backend} and are
labelled \texttt{agent\_backend: fallback} in every manifest. In fallback mode
the two LLM controllers are \emph{KPI-identical to each other and to
\texttt{rule\_based}}. The fallback ablation measures the verification
machinery, and under a perfect proposer that machinery is a no-op.

The arms are now asymmetric by construction (P0-3 fix): \texttt{llm\_agent}
applies proposals as proposed, while \texttt{llm\_agent\_verified} re-prompts
once with the violation text and falls back to the deterministic rule-based
action if still invalid (repair cap 2). Whether verification buys anything is
therefore a real-LLM question (Section~\ref{sec:m2a}).

\section{Reproducibility}\label{sec:repro}

Each run writes a \texttt{manifest.json} with the scenario config SHA-256, the
git commit of the executing tree, \texttt{git\_dirty}, package versions,
controller name, \texttt{agent\_backend}, \texttt{prompt\_version}/%
\texttt{prompt\_hash}, sampling temperature/seed and the random seed.
\texttt{failed\_runs.json} is written even when empty.

\section{Real-LLM Ablation: M2a Screen}\label{sec:m2a}

The M2a screen ran 45 24-step scenarios on the hosted API (model
\texttt{poolside/laguna-xs-2.1} via NVIDIA NIM). Both LLM arms report
\texttt{llm\_fallback\_step\_share} $= 0.000$, first-attempt validity $1.000$,
0 parse failures, 0 fenced outputs, 0 unrepaired violations, latency
p50 4134~ms / p95 8678~ms, and KPI deltas identical to the fallback-mode
\texttt{rule\_based} row.

\subsection{M2b Gate Verdict: Null Result}

The M2b gate requires $\geq 10$~pp divergence on
\texttt{llm\_unrepaired\_violations} or \texttt{llm\_fallback\_step\_share},
or a KPI CI that excludes zero. Neither condition is met: both arms report
0.0 fallback share and 0 unrepaired violations, with identical KPI CIs.
M2b is therefore \textbf{not warranted}; the screen is reported as a null.

\subsection{Validity Gate}

\texttt{llm\_fallback\_step\_share} $= 0.0 < 0.25$: the validity gate is
\textbf{passed} (not fallback-dominated). The 0.0 fallback rate means the
proposer never fell back to the deterministic proposer; every step executed an
LLM proposal. The absence of divergence is due to proposer competence, not
proposer silence.

\subsection{Null Ablation Discussion}

The null result has two interpretations: (1)~proposer competence ---
\texttt{poolside/laguna-xs-2.1} emits in-bounds, parseable JSON on every call;
and (2)~horizon effect --- the 24-step screen covers 1 simulated day and does
not stress the proposer enough. A future model/horizon change would reopen the
screen.

\section{Aggressive Failure Injection}\label{sec:aggressive}

Phase~3 added aggressive failure injection to Scenario~E: grid limit reduced by
95\%, PV inverter, battery, EV charger and sensor failures at severity 1.0.

\begin{table}[t]
\centering
\caption{Aggressive-config stress test (95\% grid limit reduction).}
\label{tab:aggressive}
\small
\begin{tabular}{lrrr}
\toprule
Controller & re\_dispatch & voltage\_viol. & resolution\_rate \\
\midrule
rule\_based & 167 & 8232 & 0.0 \\
mpc & 167 & 8232 & 0.0 \\
multi\_agent & 167 & 8232 & 0.0 \\
centralized & 167 & 8232 & 0.0 \\
\bottomrule
\end{tabular}
\end{table}

The 95\% grid limit reduction does not trigger import curtailment because the
community is net-exporting during the failure window. The 8,232 voltage
violations originate from the tightened voltage limits (0.97--1.03~pu), not
from the failures; they persist across all controllers because they are
structural. \texttt{resolution\_rate} $= 0.0$ is therefore correct: the
violations are never resolved because they cannot be, not because the
controller failed to attempt re-dispatch (167 events per controller).

\section{Threat Model and Limitations}\label{sec:threat}

\paragraph{Threat model.} We consider an adversary with full visibility of the
community's observations and actions, who may inject sensor faults, grid-limit
reductions, and component failures. Attacks are limited to the
failure-injection surface; the adversary cannot modify the scenario YAML,
controller code, or simulator internals.

\paragraph{Limitations.}
\begin{enumerate}
\item \textbf{Synthetic traces.} Load/PV profiles depend only on (seed,
  filename); reported scenario effects are device-mix effects.
\item \textbf{Controller scope.} The \texttt{rl} controller is excluded from
  the LMP-regime matrix (no tracked LMP policies at submission time);
  138-state LMP RL policies are available for a follow-up matrix.
\item \textbf{LLM proposer.} The real-LLM ablation uses a single model at
  temperature 0 on a 24-step horizon; the null result is specific to this
  combination.
\item \textbf{Aggressive config.} The 95\% grid limit does not produce import
  curtailment (community is net-exporting); the stress test exercises
  violation detection but not the curtailment path.
\item \textbf{Seed count.} Five seeds give narrow CIs; significance is
  assessed on paired deltas with Holm correction.
\end{enumerate}

\section{Caveats}

The following caveats must appear in the paper, in this order:
\begin{enumerate}
\item Resolved before the paper matrix (controller-action drop, schema
  non-propagation).
\item EV-completion artifact (root cause: missing EV definitions in schema).
\item Metric-collection defects on the pre-fix base (NaN cost, empty
  per-member costs).
\item Scenario fidelity (synthetic traces depend on seed/filename only).
\item Fallback-mode ablation (no claims about LLM proposal quality).
\item Seed count (narrow CIs; paired deltas with Holm correction).
\item Runtime tail (retired for the fixed matrix).
\end{enumerate}
